\abstract{
Data-driven methods promise tremendously fast solutions to PDEs in myriad domains, including disaster forecasting. However, as input dimensionality increases—such as spatially distributed properties, forcings, or geometries— empirical learning inevitably requires burgeoning datasets. Here, we discover a sensitivity-driven scaling principle: when physical sensitivities are enforced during learning, the information content per sample grows with input dimensionality to offset the increasing uncertainty. Harnessing this principle mitigates the curse of dimensionality and enables robust forward and inverse surrogate (neural-operator) inference with gridded inputs using far fewer samples. Consequently, we can reconstruct earthquake-induced seafloor deformation from sparse buoy observations and forecast tsunami propagation with solver-level physical consistency in under five minutes. Across diverse nonlinear PDEs spanning fluids, geophysics, and engineering, sensitivity learning acts as a scalable information channel that reshapes the data–compute trade-off and strengthens out-of-distribution robustness and long-horizon stability. Collectively, these findings extend high-fidelity physical inference to time-critical regimes previously deemed inaccessible.}
